\ulohahead{specializace UI-SU}{Kombinace modelu: bagging, boosting, náhodné lesy}
\begin{zadani}
\begin{enumerate}[label=(\alph*)]
  \item \bod{1} Co je kombinace více modelu (ensemble)? Vysvětlete rozdíl mezi baggingem a boostingem.
  \item \bod{2} Popište náhodný les (random forest) a vysvětlete, proč zlepšuje přesnost oproti jednomu stromu.
  \item \bod{3} Popište princip AdaBoostu. Za jakých podmínek ensemble pomáhá nejvíc a kdy vůbec ne (majoritní hlasování)?
\end{enumerate}
\end{zadani}

\begin{reseni}
\textbf{(a)} \emph{Ensemble} kombinuje predikce více modelu (hlasování/průměr) do silnějšího. \emph{Bagging} trénuje modely \emph{nezávisle} na bootstrap vzorcích a průměruje (snižuje hlavně \emph{rozptyl}). \emph{Boosting} trénuje modely \emph{postupně}, každý se zaměří na chyby předchozích (snižuje hlavně \emph{vychýlení}).

\textbf{(b)} \emph{Random forest} = bagging rozhodovacích stromu, kde se navíc v každém štěpení uvažuje jen náhodná podmnožina příznaku. To stromy \emph{dekoreluje}; protože průměr nekorelovaných chyb má menší rozptyl, les generalizuje lépe než jeden hluboký (přeučený) strom, a přitom téměř nepotřebuje ladění.

\textbf{(c)} \emph{AdaBoost}: udržuje váhy příkladu, v každém kole natrénuje slabý klasifikátor, zvýší váhy špatně klasifikovaných a přidá klasifikátor s váhou podle jeho přesnosti; finální predikce = vážené hlasování. Ensemble pomáhá nejvíc, jsou-li dílčí modely \emph{lepší než náhoda} a dělají \emph{různé/nekorelované} chyby. Naopak nepomůže, dělají-li všechny stejné chyby (dokonale korelované) - pak majoritní hlasování dopadne stejně jako jeden model; nejlepší případ je, když se chyby nepřekrývají a většina vždy ,,přehlasuje'' menšinu chybujících.
\end{reseni}

\begin{jadro}
Bagging = trénuj nezávisle na bootstrap vzorcích + průměruj (snižuje rozptyl). Boosting = sekvenčně, každý model na chybách předchozích (snižuje vychýlení). Random forest = bagging stromu + náhodná podmnožina příznaku (dekorelace). Pomáhá při nekorelovaných chybách.
\end{jadro}

\kalibrace{Kombinace modelu (~22\,\%) je vděčná konceptuální otázka (žádný výpočet). Bagging vs boosting + random forest = 2-3 body za vysvětlení.}
\past{Bagging $\neq$ boosting: bagging paralelně/nezávisle (rozptyl), boosting sekvenčně na chybách (vychýlení). Les pomáhá díky \emph{dekorelaci} stromu, ne jen jejich počtu.}
